
\paragraph{\textbf{Lemma 2.14.}}
\SGARefTarget{sga3:xx:2.14}
\emph{Under the preceding conditions, the following diagrams are
commutative:}

\begin{center}
  \begin{tikzcd}[column sep=huge,row sep=large]
    \mathrm W(\mathfrak g^\alpha)
      \arrow[r,"\exp"]
      \arrow[d,"\mathrm{can}"',"\sim"]
      & G
        \arrow[d,"\mathrm{can}"']
      & \mathrm W(\mathfrak g^{-\alpha})
        \arrow[l,"\exp"]
        \arrow[d,"\mathrm{can}"',"\sim"] \\
    \mathrm W(\mathfrak g'^{\alpha'})
      \arrow[r,"\exp"]
      & G'
      & \mathrm W(\mathfrak g'^{-\alpha'})
        \arrow[l,"\exp"]
  \end{tikzcd}

  \medskip
  \begin{tikzcd}[column sep=huge,row sep=large]
    \mathfrak g^\alpha\otimes\mathfrak g^{-\alpha}
      \arrow[r,"\sim"]
      \arrow[d,"\mathrm{can}"',"\sim"]
      & \mathcal O_S
        \arrow[d,"\mathrm{id}"] \\
    \mathfrak g'^{\alpha'}\otimes\mathfrak g'^{-\alpha'}
      \arrow[r,"\sim"]
      & \mathcal O_S
  \end{tikzcd}

  \medskip
  \begin{tikzcd}[column sep=huge,row sep=large]
      & T
        \arrow[dr,"\alpha"]
        \arrow[dd,"\mathrm{can}"] & \\
    \mathbb G_{m,S}
      \arrow[ur,"\alpha^\ast"]
      \arrow[dr,"\alpha'^\ast"']
      & & \mathbb G_{m,S} \\
      & T'
        \arrow[ur,"\alpha'"'] &
  \end{tikzcd}
\end{center}

\section{The Weyl group}
\SGARefTarget{sga3:xx:section:3}

\paragraph{\textbf{Notations 3.0.}}
\SGARefTarget{sga3:xx:3.0}
If \((G,T,\alpha)\) is an \(S\)-elementary system, write
\[
  N=\operatorname{Norm}_G(T),
  \qquad
  W=\operatorname{Norm}_G(T)/T
\]
\textup{(cf.\ Exposé XIX, 6.3)}.  The group \(N\) is a closed
subgroup of \(G\), smooth over \(S\).  Write
\[
  N^\times=N-T
\]
for the open subscheme of \(N\) induced on the complement of
\(T\).

Let \(R\) be the unique maximal torus of \(\Ker(\alpha)\), and let
\(T'\) be the image of
\[
  \alpha^\ast\colon\mathbb G_{m,S}\longrightarrow T;
\]
it is a one-dimensional subtorus of \(T\).  The morphism
\[
  T'\times_S R\longrightarrow T
\]
induced by multiplication in \(T\) is surjective, hence faithfully
flat.  Indeed, it is enough to check this on the geometric fibers,
where it follows at once from the formula
\[
  \alpha\circ\alpha^\ast=2.
\]

\paragraph{\textbf{Theorem 3.1.}}
\SGARefTarget{sga3:xx:3.1}
\emph{With the preceding notation:}

\smallskip
\noindent\emph{\textup{(i)} \(W\) is isomorphic to the constant group
\((\mathbb Z/2\mathbb Z)_S\).}

\smallskip
\noindent\emph{\textup{(ii)} \(N^\times\) is a locally trivial
principal homogeneous bundle under \(T\), on the left by the law
\((t,q)\mapsto tq\), and on the right by the law
\((q,t)\mapsto qt\).}

\smallskip
\noindent\emph{\textup{\SGAInlineRefTarget{sga3:xx:3.1:item:iii}(iii)} One has}
\[
  \intaut(w)t=t\cdot\alpha^\ast\bigl(\alpha(t)^{-1}\bigr)
\]
\emph{for \(w\in N^\times(S')\), \(t\in T(S')\), and \(S'\to S\).
In the decomposition \(T_{S'}=T'_{S'}\cdot R_{S'}\),
\(\intaut(w)\) induces the identity on \(R_{S'}\) and inversion on
\(T'_{S'}\).  One has the relations}
\[
  \alpha\circ\intaut(w)=\alpha^{-1},
  \qquad
  \intaut(w)\circ\alpha^\ast=(\alpha^\ast)^{-1}.
\]

\smallskip
\noindent\emph{\textup{\SGAInlineRefTarget{sga3:xx:3.1:item:iv}(iv)} For
\(X\in W(\mathfrak g^\alpha)^\times(S')\), \(S'\to S\), set}
\[
  w_\alpha(X)
    =\exp(X)\exp(-X^{-1})\exp(X).
\]
\emph{Then \(w_\alpha(X)\in N^\times(S')\), and the morphism
\(w_\alpha\colon W(\mathfrak g^\alpha)^\times\to N^\times\) so
defined satisfies}
\[
  w_\alpha(zX)
    =\alpha^\ast(z)w_\alpha(X)
    =w_\alpha(X)\alpha^\ast(z)^{-1}
\]
\emph{for \(z\in\mathbb G_m(S')\),
\(X\in W(\mathfrak g^\alpha)^\times(S')\), and \(S'\to S\).}

\smallskip
\noindent\emph{\textup{\SGAInlineRefTarget{sga3:xx:3.1:item:v}(v)} For
\(X,Y\in W(\mathfrak g^\alpha)^\times(S')\), \(S'\to S\), one has,
with the notation of \SGARefLink{sga3:xx:2.6}{2.6},}
\[
  w_\alpha(X)w_\alpha(Y)=\alpha^\ast(-XY^{-1}).
\]
\emph{In particular,}
\[
  w_\alpha(X)^2
    =\alpha^\ast(-1)
    \in {}_2T(S)\cap\operatorname{Centr}(G)(S),
\]
\[
  w_\alpha(X)^{-1}
    =w_\alpha(-X)
    =\alpha^\ast(-1)w_\alpha(X).
\]

\smallskip
\noindent\emph{\textup{\SGAInlineRefTarget{sga3:xx:3.1:item:vi}(vi)} If, similarly, one defines for
\(Y\in W(\mathfrak g^{-\alpha})^\times(S')\)}
\[
  w_{-\alpha}(Y)
    =\exp(Y)\exp(-Y^{-1})\exp(Y),
\]
\emph{then, in addition to the formulas analogous to the preceding
ones,}
\[
  w_{-\alpha}(X^{-1})
    =w_\alpha(X)^{-1}
    =w_\alpha(-X),
\qquad
  w_\alpha(X)w_{-\alpha}(Y)=\alpha^\ast(XY).
\]
